\section{Controller Interface and Rule Baseline}
\label{sec:controller}

A central usability goal of \mra{} is to integrate a custom autonomy stack through a small interface without modifying the runner, referee or adapter. This section specifies that interface, the dynamic loading mechanism, the action mapping and the two camera-assisted rule controllers used in the reported single-vehicle comparison.

\subsection{Controller Interface}

A controller implements three methods: \cmd{reset(mission\_info)}, called once before the run with a minimal static assignment; \cmd{step(observation)}, which receives the official observation of \eqref{eq:policy} at every control tick and returns the command of \eqref{eq:command}; and \cmd{close()}, called at teardown. The contract is duck-typed: any object exposing the three callables is accepted, so a participant need not depend on framework base classes. A manual controller may raise a dedicated stop exception from \cmd{step} to end a run cleanly; other exceptions are recorded as controller errors. Listing~\ref{lst:controller} shows a minimal example of local progression.

\begin{lstlisting}[language=Python, caption={Minimal custom controller.}, label={lst:controller}]
class MyController:
    def reset(self, mission_info):
        self.tracker = LocalCourseTracker(
            initial_beacon_id=mission_info[
                "initial_beacon_id"],
            total_beacons=mission_info[
                "total_beacons"],
            laps=mission_info["laps"])

    def step(self, observation):
        sensors = observation["sensors"]
        state = self.tracker.update(
            local_time_s=observation["local_time_s"],
            beacons=observation["beacons"],
            camera_image=sensors.get("FrontCamera"),
            dvl_velocity=sensors.get("DVLSensor"))
        if state.finished:
            return {"surge": 0.0, "sway": 0.0,
                    "heave": 0.0, "yaw": 0.0}
        # ... map state.beacon and camera to a command ...
        return {"surge": 0.3, "sway": 0.0,
                "heave": 0.0, "yaw": 0.0}

    def close(self):
        pass
\end{lstlisting}

The high-level command fields are listed in Table~\ref{tab:commands}; values are normalized to $[-1,1]$ and clamped by the framework before the action mapping. In fleet mode a controller may additionally return a small \cmd{message} payload alongside its command and read incoming teammate messages from the observation, using the optional inter-vehicle acoustic channel of Section~\ref{sec:fleet}.

\begin{table}[t]
\centering
\caption{High-level controller command fields of \eqref{eq:command}.}
\label{tab:commands}
\mratablestyle
\begin{tabular}{p{0.16\columnwidth}p{0.72\columnwidth}}
\toprule
Field & Interpretation \\
\midrule
\cmd{surge} & Longitudinal body-frame thrust; positive moves forward. \\
\cmd{sway} & Lateral body-frame thrust; sign follows the adapter mixing of \eqref{eq:thruster_mix}. \\
\cmd{heave} & Vertical thrust; additionally constrained by depth-safety logic near the bounds. \\
\cmd{yaw} & Yaw-rate command; sign follows the adapter mixing of \eqref{eq:thruster_mix}. \\
\bottomrule
\end{tabular}
\end{table}

\subsection{Dynamic Loading}

A controller is selected at run time in one of three ways: a built-in alias such as \cmd{rule\_gate\_baseline} or \cmd{rule\_gate\_center\_then\_commit}; a dotted module path with an explicit class; or a file path together with \cmd{--controller-class}, which is imported dynamically. For inspection and data collection, the runner can also be launched with the manual \cmd{keyboard} or \cmd{pygame} controllers. This lets an external policy be evaluated without modifying the package:

\begin{lstlisting}[language=bash]
python -m marine_race_arena.scripts.run_marine_race \
  --track .../marine_race_horseshoe_bay.json \
  --benchmark-task clean_gate \
  --controller path/to/my_controller.py \
  --controller-class MyController \
  --adapter holoocean --official --headless
\end{lstlisting}

\subsection{Action Mapping}

The framework clamps the returned command to $[-1,1]^4$ and applies a depth-safety rule that suppresses heave that would drive the vehicle past $z_{\min}$ or $z_{\max}$. The selected adapter then maps this portable body-frame command to its vehicle; the reference HoloOcean adapter uses \eqref{eq:thruster_mix}. That adapter may alternatively receive raw eight-thruster values, while the high-level interface remains the vehicle-independent default.

\subsection{Controller-Local Course Progression}

Both official rule controllers own a \cmd{LocalCourseTracker}; the runner supplies no progression signal. The tracker initializes from the public mission assignment \cmd{initial\_beacon\_id} and maintains a local expected beacon, completed count, lap and status through
\begin{equation*}
\begin{aligned}
\textsc{search}&\rightarrow\textsc{approach}\rightarrow\textsc{visual\_align}\\
&\rightarrow\textsc{commit}\rightarrow\textsc{verify\_exit}\rightarrow\textsc{advance},
\end{aligned}
\end{equation*}
ending in \textsc{finished} after the final locally confirmed beacon. It filters the received list for its own expected id; packets from other transmitters are not labeled or selected by the framework. Temporary visual loss, a missing packet, an isolated range jump, a range rise without a prior close approach and a commit without forward DVL displacement cannot advance the state.

The tracker enters and completes a passage only when persistent camera, acoustic and DVL evidence agrees with the expected beacon. Incomplete or inconsistent evidence returns it to local acquisition without advancing. The two reference controllers and their local-progression logic were designed and validated for the gate geometry, beacon placement, sensing configuration and vehicle dynamics of the three official circuits used in this study. After each non-final local advancement, the tracker holds an exit-clearance phase; after the final advancement, it enters \textsc{Finished} directly and commands zero. Controller-local snapshots and independent referee crossings are logged separately for offline comparison and never fed back into control.

\subsection{Rule-Based Baselines}

The two official reference controllers are deterministic and interpretable rule-based baselines rather than optimized policies. They use the same onboard inputs: received acoustic beacon packets, the forward RGB camera, depth and IMU measurements and body-frame DVL velocity. Both combine acoustic guidance, visual centering, depth hold and DVL-based passage verification through the shared \textsc{LocalCourseTracker} shown in Fig.~\ref{fig:rule_baseline_state}.

\emph{Continuous Servo} steers toward the expected beacon and uses the camera to refine alignment as the gate becomes visible. Depth feedback maintains the vertical reference, while DVL motion provides evidence that the vehicle has crossed and cleared the gate. Its defining behavior appears during \textsc{commit} and \textsc{verify\_exit}: the controller continues applying bounded visual corrections while moving through the aperture. It therefore keeps servoing toward the observed image center throughout the passage.

\begin{figure}[t]
\centering
\scriptsize
\begin{tikzpicture}[
  node distance=0.75cm,
  state/.style={draw, rounded corners=2pt, fill=mralight, align=center,
    minimum width=2.6cm, minimum height=0.6cm, inner sep=2.5pt},
  arrow/.style={-Latex, thick},
  lbl/.style={font=\tiny, fill=white, inner sep=1.5pt}
]
\node[state] (search) {\textsc{search}};
\node[state, below=of search] (approach) {\textsc{approach}};
\node[state, below=of approach] (align) {\textsc{visual\_align}};
\node[state, below=of align] (commit) {\textsc{commit}};
\node[state, below=of commit] (verify) {\textsc{verify\_exit}};
\node[state, below=of verify] (advance) {\textsc{advance}};

% forward path, straight down the centre
\draw[arrow] (search) -- node[lbl, right] {expected packet} (approach);
\draw[arrow] (approach) -- node[lbl, right] {$|\beta|$ small} (align);
\draw[arrow] (align) -- node[lbl, right] {centered frames} (commit);
\draw[arrow] (commit) -- node[lbl, right] {forward DVL} (verify);
\draw[arrow] (verify) -- node[lbl, right] {exit evidence} (advance);

% lost-view / failed-commit feedback returns to approach, routed on the right
\draw[arrow] (commit.east) -- ++(0.7,0)
  |- node[lbl, pos=0.5] {lost view} (approach.east);

% next-gate loop, routed on the left
\draw[arrow] (advance.west) -- ++(-0.7,0)
  |- node[lbl, pos=0.5] {next gate} (search.west);
\end{tikzpicture}
\caption{The \textsc{LocalCourseTracker} state machine shared by both rule controllers. \textsc{advance} loops to \textsc{search} for the next gate and reaches \textsc{finished} after the last confirmed passage.}
\label{fig:rule_baseline_state}
\end{figure}

\subsection{A Center-Then-Commit Rule Controller}
\label{sec:center_commit_baseline}

\emph{Center-then-Commit} retains the same observations, acoustic and visual guidance, shared tracker and local confirmation criterion as Continuous Servo. The only difference is the gate-passage strategy. After establishing a stable visual-acoustic lock, it stops chasing subsequent image-centroid changes during \textsc{commit} and \textsc{verify\_exit}. It instead advances through the aperture while maintaining the locked depth and attitude from onboard depth and IMU feedback; the same DVL-based evidence is used to verify the passage.

Because the observations, \textsc{LocalCourseTracker} and confirmation logic remain unchanged, the matched comparison of Section~\ref{sec:eval_comparison} isolates the effect of continuous visual correction versus a stabilized committed passage.

\begin{figure}[t]
\centering
\scriptsize
\begin{tikzpicture}[
  node distance=0.62cm,
  state/.style={draw, rounded corners=2pt, fill=mralight, align=center,
    minimum width=2.35cm, minimum height=0.55cm, inner sep=2.5pt},
  arrow/.style={-Latex, thick},
  lbl/.style={font=\tiny, fill=white, inner sep=1.2pt}
]
\node[state] (center) {Visual centering};
\node[state, below=of center] (commit) {Commit passage};
\node[state, below=of commit] (exit) {Exit hold};

\draw[arrow] (center) -- node[lbl, right] {four centered frames} (commit);
\draw[arrow] (commit) -- node[lbl, right] {valid crossing} (exit);
\draw[arrow] (commit.west) -- ++(-0.75,0)
  |- node[lbl, pos=0.48] {300-step timeout} (center.west);
\end{tikzpicture}
\caption{Center-then-commit gate passage and timeout recovery.}
\label{fig:center_commit_states}
\end{figure}



